%% Created by Maple 2022.1, Windows 10
%% Source Worksheet: lab_sols09
%% Generated: Wed May 08 10:35:32 BST 2024
\documentclass{article}
\usepackage{amssymb}
\usepackage{graphicx}
\usepackage{hyperref}
\usepackage{listings}
\usepackage{mathtools}
\usepackage{maple}
\usepackage[utf8]{inputenc}
\usepackage[svgnames]{xcolor}
\usepackage{amsmath}
\usepackage{breqn}
\usepackage{textcomp}
\begin{document}
\lstset{basicstyle=\ttfamily,breaklines=true,columns=flexible}
\pagestyle{empty}
\DefineParaStyle{Maple Bullet Item}
\DefineParaStyle{Maple Heading 1}
\DefineParaStyle{Maple Warning}
\DefineParaStyle{Maple Heading 4}
\DefineParaStyle{Maple Heading 2}
\DefineParaStyle{Maple Heading 3}
\DefineParaStyle{Maple Dash Item}
\DefineParaStyle{Maple Error}
\DefineParaStyle{Maple Title}
\DefineParaStyle{Maple Ordered List 1}
\DefineParaStyle{Maple Text Output}
\DefineParaStyle{Maple Ordered List 2}
\DefineParaStyle{Maple Ordered List 3}
\DefineParaStyle{Maple Normal}
\DefineParaStyle{Maple Ordered List 4}
\DefineParaStyle{Maple Ordered List 5}
\DefineCharStyle{Maple 2D Output}
\DefineCharStyle{Maple 2D Input}
\DefineCharStyle{Maple Maple Input}
\DefineCharStyle{Maple 2D Math}
\DefineCharStyle{Maple Hyperlink}
\begin{center}
\begin{Maple Title}
\underline{\textbf{Maple revision}}
\end{Maple Title}
\end{center}
\begin{Maple Normal}
\_\_\_\_\_\_\_\_\_\_\_\_\_\_\_\_\_\_\_\_\_\_\_\_\_\_\_\_\_\_\_\_\_\_\_\_\_\_\_\_\_\_\_\_\_\_\_\_\_\_\_\_\_\_\_\_\_\_\_\_\_\_\_\_\_\_\_\_\_\_\_\_\_\_\_\_\_\_\_\_\_
\end{Maple Normal}
\begin{Maple Heading 2}
\textbf{Exercise 1}
\end{Maple Heading 2}
\begin{lstlisting}
> x[1] := sin(theta) * sin(phi);
\end{lstlisting}
\begin{lstlisting}
> x[2] := sin(theta) * cos(phi);
\end{lstlisting}
\begin{lstlisting}
> x[3] := cos(theta);
\end{lstlisting}
\begin{lstlisting}
> simplify(x[1]^2 + x[2]^2 + x[3]^2);
\end{lstlisting}
\begin{Maple Normal}
\_\_\_\_\_\_\_\_\_\_\_\_\_\_\_\_\_\_\_\_\_\_\_\_\_\_\_\_\_\_\_\_\_\_\_\_\_\_\_\_\_\_\_\_\_\_\_\_\_\_\_\_\_\_\_\_\_\_\_\_\_\_\_\_\_\_\_\_\_\_\_\_\_\_\_\_\_\_\_\_\_
\end{Maple Normal}
\begin{Maple Heading 2}
\textbf{Exercise 2}
\end{Maple Heading 2}
\begin{lstlisting}
> u := 8-(x-y)^2*(x+y)^2*(16-4*x^2-y^4);
\end{lstlisting}
\begin{lstlisting}
> v := subs( x = 2*cos(t), y = 2*sin(t), u);
\end{lstlisting}
\begin{lstlisting}
> simplify(v);
\end{lstlisting}
\begin{lstlisting}
> combine(v);
\end{lstlisting}
\begin{Maple Normal}
\_\_\_\_\_\_\_\_\_\_\_\_\_\_\_\_\_\_\_\_\_\_\_\_\_\_\_\_\_\_\_\_\_\_\_\_\_\_\_\_\_\_\_\_\_\_\_\_\_\_\_\_\_\_\_\_\_\_\_\_\_\_\_\_\_\_\_\_\_\_\_\_\_\_\_\_\_\_\_\_\_
\end{Maple Normal}
\begin{Maple Heading 2}
\textbf{Exercise 3}
\end{Maple Heading 2}
\begin{lstlisting}
> solve(\{a*x+b*y+c*z=1,
       a*y+b*z+c*x=1,
       a*z+b*x+c*y=1\},
      \{x,y,z\});
\end{lstlisting}
\begin{Maple Normal}
\_\_\_\_\_\_\_\_\_\_\_\_\_\_\_\_\_\_\_\_\_\_\_\_\_\_\_\_\_\_\_\_\_\_\_\_\_\_\_\_\_\_\_\_\_\_\_\_\_\_\_\_\_\_\_\_\_\_\_\_\_\_\_\_\_\_\_\_\_\_\_\_\_\_\_\_\_\_\_\_\_
\end{Maple Normal}
\begin{Maple Heading 2}
\textbf{Exercise 4}
\end{Maple Heading 2}
\begin{lstlisting}
> fsolve(x^4+sin(x)=10^4,\{x=10\});
\end{lstlisting}
\begin{Maple Normal}
\_\_\_\_\_\_\_\_\_\_\_\_\_\_\_\_\_\_\_\_\_\_\_\_\_\_\_\_\_\_\_\_\_\_\_\_\_\_\_\_\_\_\_\_\_\_\_\_\_\_\_\_\_\_\_\_\_\_\_\_\_\_\_\_\_\_\_\_\_\_\_\_\_\_\_\_\_\_\_\_\_
\end{Maple Normal}
\begin{Maple Heading 2}
\textbf{Exercise 5}
\end{Maple Heading 2}
\begin{lstlisting}
> f := (x) -> (exp(x) + x^7)/(exp(x) - x^7);
\end{lstlisting}
\begin{lstlisting}
> evalf(f(5));
\end{lstlisting}
\begin{lstlisting}
> seq(  evalf(f(n)),  n = 0.. 50);
\end{lstlisting}
\begin{lstlisting}
> 
\end{lstlisting}
\begin{Maple Normal}
\_\_\_\_\_\_\_\_\_\_\_\_\_\_\_\_\_\_\_\_\_\_\_\_\_\_\_\_\_\_\_\_\_\_\_\_\_\_\_\_\_\_\_\_\_\_\_\_\_\_\_\_\_\_\_\_\_\_\_\_\_\_\_\_\_\_\_\_\_\_\_\_\_\_\_\_\_\_\_\_\_
\end{Maple Normal}
\begin{Maple Heading 2}
\textbf{Exercise 6}
\end{Maple Heading 2}
\begin{lstlisting}
> with(plots):
\end{lstlisting}
\begin{lstlisting}
> display(
 implicitplot(x^4+y^4=4,x=-3..3,y=-3..3),
 plot([(2+sin(8*t))*cos(t),(2+sin(8*t))*sin(t),t=0..2*Pi])
);
\end{lstlisting}
\begin{Maple Normal}
\_\_\_\_\_\_\_\_\_\_\_\_\_\_\_\_\_\_\_\_\_\_\_\_\_\_\_\_\_\_\_\_\_\_\_\_\_\_\_\_\_\_\_\_\_\_\_\_\_\_\_\_\_\_\_\_\_\_\_\_\_\_\_\_\_\_\_\_\_\_\_\_\_\_\_\_\_\_\_\_\_
\end{Maple Normal}
\begin{Maple Heading 2}
\textbf{Exercise 7}
\end{Maple Heading 2}
\begin{lstlisting}
> plot(
[tan(Pi*x),cot(Pi*x),-1,1],    # functions to plot
  x=-4..4,                     # horizontal range
    -2..2,                     # vertical range
  scaling=constrained,         # same scale on both axes
  discont=true                 # skip over discontinuities
);
\end{lstlisting}
\begin{Maple Normal}
\_\_\_\_\_\_\_\_\_\_\_\_\_\_\_\_\_\_\_\_\_\_\_\_\_\_\_\_\_\_\_\_\_\_\_\_\_\_\_\_\_\_\_\_\_\_\_\_\_\_\_\_\_\_\_\_\_\_\_\_\_\_\_\_\_\_\_\_\_\_\_\_\_\_\_\_\_\_\_\_\_
\end{Maple Normal}
\begin{Maple Heading 2}
\textbf{Exercise 8}
\end{Maple Heading 2}
\begin{lstlisting}
> restart;
\end{lstlisting}
\begin{lstlisting}
> f := (t) -> ((2+sqrt(3))*t-1)/(2+sqrt(3)+t);
\end{lstlisting}
\begin{lstlisting}
> g := (t) -> f(f(f(t)));
\end{lstlisting}
\begin{lstlisting}
> h := (t) -> g(g(t));
\end{lstlisting}
\begin{lstlisting}
> simplify(g(t));
\end{lstlisting}
\begin{lstlisting}
> simplify(h(t));
\end{lstlisting}
\begin{Maple Normal}
\_\_\_\_\_\_\_\_\_\_\_\_\_\_\_\_\_\_\_\_\_\_\_\_\_\_\_\_\_\_\_\_\_\_\_\_\_\_\_\_\_\_\_\_\_\_\_\_\_\_\_\_\_\_\_\_\_\_\_\_\_\_\_\_\_\_\_\_\_\_\_\_\_\_\_\_\_\_\_\_\_
\end{Maple Normal}
\begin{Maple Heading 2}
\textbf{Exercise 9}
\end{Maple Heading 2}
\begin{lstlisting}
> y := sin(10*x);
\end{lstlisting}
\begin{lstlisting}
> (diff(y,x$8)/y)^(1/8);
\end{lstlisting}
\begin{lstlisting}
> simplify(%);
\end{lstlisting}
\begin{lstlisting}
> 
\end{lstlisting}
\begin{Maple Normal}
\_\_\_\_\_\_\_\_\_\_\_\_\_\_\_\_\_\_\_\_\_\_\_\_\_\_\_\_\_\_\_\_\_\_\_\_\_\_\_\_\_\_\_\_\_\_\_\_\_\_\_\_\_\_\_\_\_\_\_\_\_\_\_\_\_\_\_\_\_\_\_\_\_\_\_\_\_\_\_\_\_
\end{Maple Normal}
\begin{Maple Heading 2}
\textbf{Exercise 10}
\end{Maple Heading 2}
\begin{Maple Normal}
Here is the efficient method using the \textcolor[RGB]{255,0,0}{\textbf{seq()}} command:\


\end{Maple Normal}
\begin{lstlisting}
> plot([seq(n^(3/2) * x^n * (1-x)^2,n=1..10)],x=0..1);
\end{lstlisting}
\begin{Maple Normal}
\

Here is the less efficient method with more typing:
\end{Maple Normal}
\begin{lstlisting}
> plot([
x * (1-x)^2,
2^(3/2) * x^2 * (1-x)^2,
3^(3/2) * x^3 * (1-x)^2,
4^(3/2) * x^4 * (1-x)^2,
5^(3/2) * x^5 * (1-x)^2,
6^(3/2) * x^6 * (1-x)^2,
7^(3/2) * x^7 * (1-x)^2,
8^(3/2) * x^8 * (1-x)^2,
9^(3/2) * x^9 * (1-x)^2,
10^(3/2) * x^10 * (1-x)^2
],x=0..1);
\end{lstlisting}
\begin{Maple Normal}
\_\_\_\_\_\_\_\_\_\_\_\_\_\_\_\_\_\_\_\_\_\_\_\_\_\_\_\_\_\_\_\_\_\_\_\_\_\_\_\_\_\_\_\_\_\_\_\_\_\_\_\_\_\_\_\_\_\_\_\_\_\_\_\_\_\_\_\_\_\_\_\_\_\_\_\_\_\_\_\_\_
\end{Maple Normal}
\begin{Maple Heading 2}
\textbf{Exercise 11}
\end{Maple Heading 2}
\begin{lstlisting}
> y := x^4*(1-x)^4/(1+x^2);
\end{lstlisting}
\begin{lstlisting}
> int(y,x);
\end{lstlisting}
\begin{lstlisting}
> int(y,x=0..1);
\end{lstlisting}
\begin{lstlisting}
> simplify(diff(y,x));
\end{lstlisting}
\begin{lstlisting}
> plot(y,x=0..1);
\end{lstlisting}
\begin{Maple Normal}
The hump has position 
$p$ and height 
$h$, where 
$p$ is the root of 
$\frac{\mathit{dy}}{\mathit{dx}}=0$ near 
$x = 0.5$ and 
$h$ is obtained\

by putting 
$x =p$ in 
$y$.
\end{Maple Normal}
\begin{lstlisting}
> p := fsolve(diff(y,x)=0,x=0.5);
\end{lstlisting}
\begin{lstlisting}
> h := subs(x=p, y);
\end{lstlisting}
\begin{Maple Normal}
\_\_\_\_\_\_\_\_\_\_\_\_\_\_\_\_\_\_\_\_\_\_\_\_\_\_\_\_\_\_\_\_\_\_\_\_\_\_\_\_\_\_\_\_\_\_\_\_\_\_\_\_\_\_\_\_\_\_\_\_\_\_\_\_\_\_\_\_\_\_\_\_\_\_\_\_\_\_\_\_\_
\end{Maple Normal}
\begin{Maple Heading 2}
\textbf{Exercise 12}
\end{Maple Heading 2}
\begin{lstlisting}
> a := (n) -> evalf((1+1/n)^n/exp(1));
\end{lstlisting}
\begin{lstlisting}
> seq(a(n),n=1..100);
\end{lstlisting}
\begin{Maple Normal}
The first number greater than 0.99 occurs around the middle of the list, close to 
$n =50$.  We therefore\

look at a shorter list centred around that value:
\end{Maple Normal}
\begin{lstlisting}
> seq(a(n),n=48..52);
\end{lstlisting}
\begin{Maple Normal}
We see that 
$a (50)= 0.9901799000$, and that this is the first number greater than 0.99.\


\end{Maple Normal}
\begin{Maple Normal}
Here is a slicker way:
\end{Maple Normal}
\begin{lstlisting}
> min(op(select(n -> (a(n)>0.99),[seq(n,n=1..100)])));
\end{lstlisting}
\begin{lstlisting}
> 
\end{lstlisting}
\begin{lstlisting}
> 
\end{lstlisting}
\begin{lstlisting}
> 
\end{lstlisting}
\begin{lstlisting}
> 
\end{lstlisting}
\end{document}
